\PassOptionsToPackage{log-declarations=false}{xparse}
% \documentclass{article}
% \usepackage{fontspec}
% \setmainfont{FandolSong}

% \begin{document}
% Hello, world!
% \end{document}

\documentclass[twoside,nofonts]{fduthesis}
\usepackage{kantlipsum}
\usepackage{zhlipsum}
% \ExplSyntaxOn
% \bool_show:N\g__fdu_no_font_bool
% \ExplSyntaxOff

\fdusetup{
  style = {
    % font = libertinus,
    % cjkfont = founder,
    fullwidthstop,
    footnotestyle = plain,
    % automakecover = false
  },
  info = {
    title = {自河南经乱关内阻饥兄弟离散各在一处},
    title* = {Measurements of the interaction between energetic
      photons and hadrons show that the interaction},
%    date = {2017年2月10日},
    author = {某某某},
    author* = {Xiangdong Zeng},
    supervisor = {陈丙丁 \quad 教授},
    instructors = {
      张五六 \quad 工程师,
      赵\quad 甲 \quad 工程师,
      王三四 \quad 讲\quad 师
    },
    major = {物理学},
    department = {凝聚态物理系},
    secretlevel = ii,
    secretyear = {三年},
    studentid = {14307110000},
    keywords = {\LaTeX, hih, vwvs, 物理, 中心法则, vsbd, 为侨服务},
    keywords* = {\LaTeX, hih, vwvs, Physics, Central Rule, vsbd},
    clc = {O414.1/65}
  }
}

\setmainfont{XITS}
\setsansfont{TeX Gyre Heros}
\setmonofont{TeX Gyre Cursor}%[Ligatures=NoCommon]
\setmathfont{XITS Math}

\setCJKmainfont[ItalicFont=FandolFang]{FandolSong}
\setCJKsansfont{FandolHei}

\newcommand{\fonttesttext}{你好，世界\symbol{12290}Hello, world!}
\newcommand\fonttest{%
  正常：\qquad\fonttesttext \par
  粗体：\qquad\textbf{\fonttesttext} \par
  倾斜：\qquad\textit{\fonttesttext} \par
  小型大写：\qquad\textsc{\fonttesttext}
}

\begin{document}
\frontmatter
\tableofcontents

\begin{abstract}
\ExplSyntaxOn
% \cs_show:N \fdu_cjk_font_kai:
\ExplSyntaxOff
\emph{中心法则}形成了分子生物学的生命观：生命世界在信息和规律上是统一的；
生命的物质基础和主宰物质是不同的，生命的物质基础是以核酸蛋白质整合
体系为主宰的原生质各种必要的物质组分及其实在的相互作用；生命运动
的本质在于组成生命的各物质之间以及生命的物质、能量与信息之间和生命
自身与环境之间的不断的相互作用。不少人不少人，如果是到无穷大。
\end{abstract}

\begin{abstract*}
\LaTeX3 does not use @ as a ``letter'' for defining internal
macros. Instead, the symbols are used in internal macro names to
provide structure. The name of each function is divided into
logical units using separates the name of the function from the
argument specifier (``arg-spec''). This describes the arguments
expected by the function.
\end{abstract*}

\mainmatter
\chapter{文本，字体，脚注 \quad Text, font and footnote}
\section{文字与段落 Text and paragraph}
\subsection{中文文本 Chinese}
\zhlipsum

\subsection{英文文本 English}
\kant

\clearpage

\section{字体 Font}
\subsection{普通字体 Roman}
\fonttest

\subsection{无衬线字体 Sans-serif}
\textsf{\fonttest}

\subsection{打字机字体 Typewriter}
\texttt{\fonttest}

\subsection{句号}
如果\symbol{"002E}会突然：\textsc{Full Stop} \par
如果\symbol{"3002}会突然：\textsc{Ideographic Full SStop} \par
如果\symbol{"FF0E}会突然：\textsc{Fullwidth Full Stop} \par
如果\symbol{"FF61}会突然：\textsc{Halfwidth Ideographic Full Stop} \par

\clearpage

\section{脚注 Footnote}
脚注\footnote{脚注1如果会突然}。

脚注 \footnote{脚注2。千千万的。}。未取得的

脚注。\footnote{脚注3是一个长脚注。\zhlipsum*[2]}未取得的

\textit{脚注倾斜。 \footnote{脚注4}未取得的}

\textbf{脚注加粗。\footnote{脚注5} 未取得的}

vs\footnote{脚注6}未取得的

vs \footnote{脚注7要分段。\par 不舒服不得不运河滩上野跑，头顶着毒热的阳光，身上再裹起兜肚，一不风凉，
二又窝汗，穿不了一天，就得起大半身痱子。再有，全村跟他一般大的小姑娘，
谁的兜肚也没有这么花儿草儿的鲜艳，他穿在身上，男不男，女不女，
小姑娘们要用手指刮破脸蛋儿。}未取得的

vs\footnote{脚注8}ge

vs\footnote{脚注9} jsty

vs\footnote{脚注10分三段。青大娘大高个儿，一双大脚，青铜肤色，
嗓门也亮堂，骂起人来，方圆二三十里，敢说找不出能够招架几个回合的敌手。
一丈青大娘骂人，就像雨打芭蕉，长短句。\par
青大娘大高个儿，一双大脚，青铜肤色，
嗓门也亮堂，骂起人来，方圆二三十里，敢说找不出能够招架几个回合的敌手。
一丈青大娘骂人，就像雨打芭蕉，长短句。 \par
青大娘大高个儿，一双大脚，青铜肤色，
嗓门也亮堂，骂起人来，方圆二三十里，敢说找不出能够招架几个回合的敌手。
一丈青大娘骂人，就像雨打芭蕉，长短句。}

\chapter{数学与定理}
\section{数学 Math}
\[\pi=\sb{33}\]

\[
  \int\sin x\,\mathrm{d}x=\cos x + C
\]

\begin{proof}
道千乘之国，敬事而信，节用而爱人，使民以时。
\end{proof}

\begin{definition}
证明完毕/证讫，又写作Q.E.D.。这是拉丁词组“quod erat demonstrandum”（这就是所要证明的）的缩写，译自希腊语“ὅπερ ἔδει δεῖξαι”（hoper edei deixai），很多早期数学家用过，包括欧几里得和阿基米德。“Q.E.D.”可以在证明的尾段写出，以显示证明所需的结论已经完整了。
\end{definition}

\begin{lemma}
  这是一条华丽丽的引理。
\end{lemma}

\begin{proof}
先帝创业未半而中道崩殂，今天下三分，益州疲弊，此诚危急存亡之秋也。
\begin{equation}
  \sum_{k=0}^{\infty} \frac{1}{x^k} = \int \sin x dx
\end{equation}
\end{proof}

\begin{proof}
先帝创业未半而中道崩殂，今天下三分，益州疲弊，此诚危急存亡之秋也。
\begin{equation*}
  \sum_{k=0}^{\infty} \frac{1}{x^k} = \int \sin x dx
\end{equation*}
\end{proof}

\begin{lemma}
  这又是一条华丽丽的引理。
\end{lemma}

% \chapter{图表 vs 浮动体}
% \section{title}
% Myriad，英语单词，意为「无数的」。同时，「Myriad」也是一款字体的名字。
% 由罗伯特·斯林巴赫（Robert Slimbach，1956年-）和卡罗·图温布利
% （Carol Twombly，1959年-）在1990年到1992年期间以 Frutiger 字体为蓝本
% 为 Adobe 公司设计。 Myriad 是早期数码字体时代的先驱，伴随着技术的成长
% 一路走来。

% \begin{figure}[h]
%   \centering
%   \includegraphics[width=3cm]{../logo/fudan-emblem.pdf}
%   \includegraphics[width=4cm]{../logo/fudan-emblem.pdf}
%   \caption{Multiple Master 是 Type 1字体格式的扩展部分。Type 1 是利用
%     PostScript 语言描述字形信息的字体系统。Type 1字体是第一款矢量字体
%     （outline font），通过二维坐标系中的关键点和三次贝塞尔曲线描述字体
%     的边缘，在屏幕显示和输出时，在光栅图像处理器内，根据字号大小计算
%     出字体边缘（栅格化）。}
% \end{figure}

% 如今，它更多地和我们相见在显示屏幕上。当然，还有那著名的标榜设计的
% 电子品牌。1992 年，耗时两年开发的 Myriad 终于发布了历史上第一个版本：
% Myriad MM。

% \section{title}
% 这款温和且具有良好可读性的人文主义无衬线字体，集诸多当时最新的数字
% 字体技术于一身。 后缀 MM，意为 Multiple Master，没有找到对应的中文
% 译名，我们权且称之为「多母板技术」。Myriad 是最早采用 Multiple Master
% 技术的无衬线字体之一。这项技术的原理是在坐标轴（Axis）的区间两端设计
% 极限母板，中间的变量则采取线性或非线性变化，对于字体来说，字型的宽度、
% 粗细甚至有无衬线，都可以在坐标轴上设置。此外，MM 技术还提供了在小字号
% 下屏幕显示的视觉修正（Optical Adjustment），也就是说，同一款字体，在
% 小字号时，其字间距和笔画粗细，会被适当地放大。而衬线字体，随着字号的
% 变小，衬线会相对变粗。视觉修正可以提高小字号字体的识别性，对于远低于
% 印刷分辨率的电脑屏幕来说，也具有重要意义。

% \begin{table}[h]
%   \centering
%   \caption{一个 normal 表格}
%   \begin{tabular}{ccc}
%     \hline
%     \bfseries 功能 & \bfseries 环境 & \bfseries code \\
%     \hline
%     表格 & tabular & \ttfamily \backslash begin\{tabular\} ... \backslash end\{tabular\} \\
%     插图 & figure  & \ttfamily \backslash begin\{figure\}  ... \backslash end\{figure\}  \\
%     居中 & center  & \ttfamily \backslash begin\{center\}  ... \backslash end\{center\}  \\
%     \hline
%   \end{tabular}
% \end{table}

% 在 Multiple Master 的时代，字号是从6pt到72pt之间非线性设置的。这一传统
% 保留到了今天 Truetype 和 Opentype 的 Single Master 时代。Adobe 软件的
% 字体下拉菜单，仍然只显示6到72pt 的字号。

\end{document}
